
seul polyn\^ome de Hilbert\,; donc (iii) r\'esulte de 2.11.

\textbf{Remarque 4.8.} Soit @@M0@@ propre et int\`egre sur @@M1@@ noeth\'erien. Sans hypoth\`ese de projectivit\'e sur @@M2@@, on ignore s'il reste vrai que @@M3@@ est une immersion \uline{ferm\'ee}, sauf lorsque @@M4@@ est normal sur @@M5@@ et @@M6@@ est \uline{repr\'esentable} (FGA 236, 2.3).

% Added in pages 654--669
\providecommand{\NS}{\operatorname{NS}}
\providecommand{\Picard}{\operatorname{Pic}}
\providecommand{\Et}{\mathrm{et}}
\providecommand{\rk}{\operatorname{rk}}
\providecommand{\im}{\operatorname{Im}}
\providecommand{\HH}{\mathrm H}


\subsection*{5. Th\'eor\`emes de finitude du type ``N\'eron-S\'everi''.}

\textbf{Th\'eor\`eme 5.1.} \emph{Soit @@M7@@ propre sur @@M8@@ noeth\'erien\,; alors les groupes de N\'eron-S\'everi @@M9@@ des fibres g\'eom\'etriques @@M10@@ de @@M11@@ sont de type fini, et leur rang et l'ordre de leur sous-groupe de torsion @@M12@@ restent major\'es.}

En effet, par r\'ecurrence noeth\'erienne, on peut se borner \`a prouver 5.1 en y rempla\c{c}ant @@M13@@ par un ouvert non-vide arbitraire de @@M14@@\,; ainsi on peut supposer @@M15@@ repr\'esentable par un @@M16@@-sch\'ema de type fini (XII 1.2 et 4.7 (iii)). Alors l'ordre de @@M17@@ est \'egal au nombre de composantes connexes de @@M18@@, nombre qui reste le m\^eme au-dessus d'un ouvert non-vide de @@M19@@ (EGA IV 9.7.9).

Il reste \`a majorer le rang de @@M20@@. Comme dans la d\'emonstration de 4.7 (iii), on se ram\`ene facilement au cas @@M21@@ int\`egre, et @@M22@@ projectif et int\`egre sur @@M23@@. Soit @@M24@@ la dimension commune des fibres de @@M25@@. Si @@M26@@ est @@M27@@ (resp. @@M28@@), alors ce rang est \'evidemment @@M29@@ (resp. @@M30@@).

Si @@M31@@ est @@M32@@, on va se ramener au cas @@M33@@ et @@M34@@ lisse sur @@M35@@. Quitte \`a restreindre @@M36@@ on peut trouver une section lin\'eaire @@M37@@ de @@M38@@ qui est
